\section{対象データと用語}
\label{sec_data}
\subsection{ボケてとお題画像}
ボケてでは，投稿者が1枚の画像を掲示し，その画像を題目とする大喜利の場を作る．
本稿では，この画像をお題画像と呼び，お題画像とその投稿の単位をお題と呼ぶ．
他の利用者がお題に対して投稿する面白い一言をボケと呼ぶ．
各お題は数値の識別子（お題 ID）を持ち，1件のお題には0件以上のボケが付く．

本稿で用いるデータは，ボケてから収集した1次データからなるデータベース BoketeDBv4 である．
BoketeDBv4 は，お題，ボケ，お題画像，利用者，評価などの表からなり，各表は Apache Parquet 形式で保存されている．
お題画像の表は6,823,270件の画像を含み，各画像のファイルの内容から求めた SHA-256 のハッシュ値を持つ．
ハッシュ値が同じ画像は内容が同一であり，同一の内容をまとめると6,609,511枚となる．
お題の投稿年は2008年から2026年にわたり，2014年から2023年に毎年40万件以上が投稿されている．

お題の表には，投稿者が投稿時に選んだ画像の種類を表す \texttt{type} 列がある．
\texttt{type} 列は2022年以降の約80万件のお題にのみ値があり，値は「写真」「イラスト」「穴埋め」「文字」「その他」のいずれかである．
また，お題には投稿者や閲覧者が付けたタグがあり，「穴埋め」「かるた」「漫画」「アニメ」などのタグを含む．
これらのメタデータは画像の種類の手がかりになるが，\ref{sec_selection}章で示すとおり，ラベルとして用いるには誤りが多い．

\subsection{お題画像の種類と処理の要件}
お題画像の種類のうち，ボケの生成の処理に影響するものは次のとおりである．
現実画像は，現実の写真であり，画像の内容を見て一言を付ければよい．
アニメのスクリーンショットと漫画のコマの切り抜きは，その画像のキャラクターや作品の知識を前提とすることがある．
文章のお題が書かれた画像は，画像中の文字を読み取ったうえで，その文章に答える必要がある．
穴埋め画像は，画像中に矩形や円などの空欄の記号があり，利用者はその空欄に入る語をボケとして投稿する．
例えば「東京大学○○学部」の画像では，○○に入る語を答える．
カルタ画像は，画像中の円などの記号の中にひらがな1文字が書かれており，利用者はそのひらがなから始まる文を投稿する．
穴埋めとカルタでは，記号の位置やひらがなを検出したうえでボケを生成する必要がある．

これらの種類を約682万件のお題画像に付与する処理には，次の要件がある．
第1に，有料の API を用いず，ローカルの計算機で処理する．
第2に，全件の処理を現実的な時間，具体的には単一の GPU で1日程度以内に終える．
第3に，ラベルの定義を見直した場合に，全件の処理をやり直す費用が小さい．
本研究の実験は全て，VRAM が96GB の NVIDIA RTX PRO 6000 Blackwell を1基備えた計算機で行った．

\subsection{本稿で用いる用語}
本稿では，ラベルの構成や処理の構成を表すために，いくつかの用語を独自の意味で用いる．
表\ref{tab_terms}に，主な用語とその意味をまとめる．
各用語は，初めて用いる箇所でも改めて説明する．

\begin{table*}[t]
    \centering
    \caption{本稿で用いる主な用語}
    \label{tab_terms}
    \small
    \begin{tabular}{>{\raggedright\arraybackslash}p{0.2\textwidth}>{\raggedright\arraybackslash}p{0.74\textwidth}}
        \toprule
        用語 & 意味 \\
        \midrule
        形式 & 大喜利の回答の形を制限するお題画像の性質．カルタ，穴埋め，お題文，通常の4クラスから1つを選ぶ \\
        お題文 & 投稿者が画像に書き加えた，回答者に向けた問い，指示，題目の句などの文．形式のクラス名としても用いる \\
        属性 & 形式とは独立に，画像中の記号や文字の有無を表す複数選択のラベル．最終の定義では空欄あり，カルタ記号あり，文字ありの3項目 \\
        媒体 & 画像の画風や由来を表すラベル．写真，アニメ，ゲーム，マンガ，テキストグラフィック，絵画，イラストの7クラスから原則1つを選ぶ \\
        軸 & 形式，属性，媒体のように，独立に判定するラベルのまとまり \\
        定義書 & ラベルの定義を文章でまとめた文書．第1版から第4版まで改訂した \\
        教師 & 学習用のラベルを付ける VLM．本稿では Qwen3-VL-8B に複数の質問をする構成全体を指す \\
        教師ラベル & 教師が画像に付けたラベル．生徒の学習に用いる \\
        生徒 & 教師ラベルで学習し，全件の処理に用いる高速な画像分類器 \\
        蒸留 & 教師の判断を，教師ラベルを用いた学習によって生徒に移すこと \\
        目視ラベル & 判断者が画像を見て付けたラベル．評価にのみ用いる \\
        通常の評価用データ & 目視ラベルを付けた640枚の評価用の画像．調整用（dev）と最終評価用（test）に二分する \\
        少ない区分 & 出現率が数\%以下のクラス．具体的にはゲーム，CG，絵画，カルタ \\
        少ない区分の評価用データ & 少ない区分を多く含むように抽出し，目視ラベルを付けた800枚の評価用の画像 \\
        実データの目視ラベル & 学習にも評価用データにも用いていない画像から抽出し，目視ラベルを付けた380枚．無作為抽出の160枚と層別抽出の220枚からなる \\
        少ない区分の追加 & 生徒の予測で少ない区分らしい画像を選び，教師ラベルを付けて学習データに加えること \\
        確認の質問 & 生徒または教師がお題文と判定した画像に限り，VLM に別の短い質問をして判定を確かめる処理 \\
        書き出しの集中度 & お題に付いたボケの最初の文字が，最も多い1文字に集中している割合．カルタの判定の補助に用いる \\
        全項目一致 & 形式，属性，媒体の全ての項目が目視ラベルと一致した画像の割合 \\
        \bottomrule
    \end{tabular}
\end{table*}
